\documentclass[10pt,a4paper]{article}
\usepackage[utf8]{inputenc}
\usepackage[T1]{fontenc}
\usepackage[margin=1.2cm,top=1.0cm,bottom=1.0cm]{geometry}
\usepackage{multicol}
\usepackage{xcolor}
\usepackage{enumitem}
\usepackage{titlesec}
\usepackage{tabularx}

% Cores
\definecolor{accent}{HTML}{2E86C1}
\definecolor{keybg}{HTML}{E8E8E8}
\definecolor{darktext}{HTML}{1A1A1A}

% Sem numera\c{c}\~ao de p\'agina
\pagestyle{empty}

% Formato de se\c{c}\~ao: sem risco, compacto
\titleformat{\section}
  {\large\bfseries\color{accent}}
  {}{0em}{}
\titlespacing*{\section}{0pt}{6pt}{3pt}

% Comando para tecla formatada
\newcommand{\key}[1]{{\small\texttt{#1}}}
\newcommand{\desc}[1]{{\small #1}}

% Comando para linha de atalho
\newcommand{\K}[2]{\key{#1} \hfill \desc{#2}\\[1.5pt]}

% Espa\c{c}amento entre colunas
\setlength{\columnsep}{14pt}
\setlength{\parindent}{0pt}

\begin{document}

\begin{center}
{\LARGE\bfseries\color{darktext} Neovim para Escrita Acad\^emica}\\[3pt]
{\small\color{gray} Refer\^encia r\'apida $\cdot$ Navegar, editar e revisar a sua tese}
\end{center}

\vspace{4pt}

\begin{multicols}{3}
\raggedcolumns

% ============================================================
\section{Comece a escrever}
\K{i}{Escrever antes do cursor}
\K{a}{Escrever depois do cursor}
\K{I}{Escrever no in\'icio da linha}
\K{A}{Escrever no final da linha}
\K{o}{Novo par\'agrafo/linha abaixo}
\K{O}{Novo par\'agrafo/linha acima}
\K{Esc}{Parar de escrever (normal)}
\desc{No modo normal as teclas comandam o texto; em inser\c{c}\~ao voc\^e digita.}

% ============================================================
\section{Navegar no texto}
\K{h / j / k / l}{Esq, baixo, cima, dir}
\K{w / b}{Pr\'oxima / anterior palavra}
\K{e}{At\'e o final da palavra}
\K{0}{In\'icio da linha}
\K{\^{}}{Primeira palavra da linha}
\K{\$}{Final da linha}
\K{( / )}{Frase anterior / pr\'oxima}
\K{\{ / \}}{Par\'agrafo anterior / pr\'oximo}
\K{gg}{In\'icio do documento}
\K{G}{Final do documento}
\K{:\{n\}}{Ir para a linha \{n\}}
\K{Ctrl+d / Ctrl+u}{Meia p\'agina baixo / cima}
\K{H / M / L}{Topo / meio / final da tela}
\K{f\{c\}}{Saltar at\'e o caractere \{c\}}
\K{; / ,}{Repetir / inverter o salto f}

% ============================================================
\section{Linhas longas (wrap)}
\desc{Quando o par\'agrafo quebra na tela:}
\vspace{3pt}
\K{gj / gk}{Descer / subir por linha visual}
\K{:set wrap}{Quebra de linha suave}
\K{:set linebreak}{N\~ao cortar palavras}

% ============================================================
\section{Editar prosa}
\K{ciw}{Trocar a palavra inteira}
\K{cis}{Trocar a frase inteira}
\K{cip}{Trocar o par\'agrafo inteiro}
\K{cc}{Reescrever a linha}
\K{c\{mov.\}}{Trocar at\'e o movimento}
\K{C}{Trocar at\'e o final da linha}
\K{daw}{Apagar a palavra (e o espa\c{c}o)}
\K{das}{Apagar a frase inteira}
\K{dap}{Apagar o par\'agrafo inteiro}
\K{dd}{Apagar a linha}
\K{x}{Apagar um caractere}
\K{r\{c\}}{Corrigir uma letra (typo)}
\K{J}{Juntar a pr\'oxima linha \`a atual}
\K{> / <}{Recuar / avan\c{c}ar indenta\c{c}\~ao}
\K{.}{Repetir a \'ultima edi\c{c}\~ao}

% ============================================================
\section{Recortar, copiar, colar}
\K{yiw / yip}{Copiar palavra / par\'agrafo}
\K{yy}{Copiar a linha}
\K{y\{mov.\}}{Copiar at\'e o movimento}
\K{dd / d\{mov.\}}{Recortar linha / movimento}
\K{p / P}{Colar depois / antes}
\K{"+y / "+p}{Copiar / colar via clipboard}

% ============================================================
\section{Mai\'uscula / min\'uscula}
\K{\~{}}{Alternar caixa da letra}
\K{guiw}{Palavra em min\'uscula}
\K{gUiw}{Palavra em MAI\'USCULA}
\K{guis / gUis}{Frase em min/MAI\'USCULA}

% ============================================================
\section{Selecionar (visual)}
\K{v}{Selecionar caractere a caractere}
\K{V}{Selecionar linha inteira}
\K{viw}{Selecionar a palavra}
\K{vis}{Selecionar a frase}
\K{vip}{Selecionar o par\'agrafo}
\desc{Depois: \texttt{d} apaga, \texttt{c} troca, \texttt{y} copia.}

% ============================================================
\section{Objetos de texto}

Use ap\'os o operador (c, d, y, v):

\texttt{i} = interno (s\'o o conte\'udo)\\
\texttt{a} = ao redor (inclui espa\c{c}o/marcas)

\vspace{3pt}
\K{iw / aw}{Palavra}
\K{is / as}{Frase}
\K{ip / ap}{Par\'agrafo}
\K{i" / a"}{Dentro / ao redor de aspas}
\K{i( / a(}{Dentro / ao redor de par\^enteses}
\K{i\{ / a\{}{Dentro / ao redor de chaves}

% ============================================================
\section{Buscar no texto}
\K{/texto}{Buscar para frente}
\K{?texto}{Buscar para tr\'as}
\K{n / N}{Pr\'oxima / anterior ocorr\^encia}
\K{*}{Buscar a palavra sob o cursor}

% ============================================================
\section{Revisar e substituir}
\K{:s/velho/novo/}{Trocar na linha (1\textsuperscript{a})}
\K{:s/velho/novo/g}{Trocar todas na linha}
\K{:\%s/velho/novo/g}{Trocar em todo o texto}
\K{:\%s/velho/novo/gc}{Trocar com confirma\c{c}\~ao}
\desc{O \texttt{c} no fim pede confirma\c{c}\~ao a cada troca: \'otimo para revis\~ao.}

% ============================================================
\section{Corretor ortogr\'afico}
\K{:set spell}{Ligar o corretor}
\K{:set spelllang=pt\_br}{Idioma portugu\^es}
\K{:set nospell}{Desligar o corretor}
\K{]s / [s}{Pr\'oximo / anterior erro}
\K{z=}{Sugest\~oes de corre\c{c}\~ao}
\K{zg}{Adicionar palavra ao dicion\'ario}
\K{zw}{Marcar palavra como incorreta}

% ============================================================
\section{Desfazer / refazer}
\K{u}{Desfazer}
\K{Ctrl+r}{Refazer}
\K{U}{Desfazer mudan\c{c}as na linha}

% ============================================================
\section{Marcas (ir e voltar)}
\K{m\{a-z\}}{Marcar este ponto do texto}
\K{`\{marca\}}{Voltar \`a marca (exata)}
\K{'\{marca\}}{Voltar \`a marca (linha)}
\K{``}{Voltar de onde voc\^e saltou}
\K{:marks}{Listar as marcas}
\desc{Pule entre cap\'itulo, notas e refer\^encias.}

% ============================================================
\section{Dobrar se\c{c}\~oes (folds)}
\desc{Recolher partes longas para ver a estrutura:}
\vspace{3pt}
\K{za}{Abrir/fechar a dobra}
\K{zR / zM}{Abrir / fechar todas}

% ============================================================
\section{Salvar e sair}
\K{:w}{Salvar}
\K{:wq / ZZ}{Salvar e sair}
\K{:q}{Sair}
\K{:q!}{Sair sem salvar}
\K{:e arquivo}{Abrir outro arquivo}
\K{:w novo.md}{Salvar com outro nome}

% ============================================================
\section{V\'arios documentos}
\desc{Cap\'itulos e notas abertos ao mesmo tempo:}
\vspace{3pt}
\K{:e arquivo}{Abrir um arquivo}
\K{:bn / :bp}{Pr\'oximo / anterior documento}
\K{:ls}{Listar documentos abertos}
\K{Ctrl+w v}{Dividir a tela ao lado}
\K{Ctrl+w s}{Dividir acima/abaixo}
\K{Ctrl+w h/j/k/l}{Pular entre as divis\~oes}

% ============================================================
\section{Combinando comandos}
\desc{A gram\'atica do Vim: verbo + quantidade + objeto.}
\vspace{3pt}
\K{ciw}{(trocar)(palavra)}
\K{das}{(apagar)(uma frase)}
\K{yap}{(copiar)(um par\'agrafo)}
\K{d2w}{(apagar)(2)(palavras)}
\K{3dd}{(apagar)(3)(linhas)}
\K{>ip}{(recuar)(par\'agrafo)}
\K{ci"}{(trocar)(dentro das aspas)}
\K{gUiw}{(MAI\'USCULA)(palavra)}

\end{multicols}
\end{document}
